% ============================================================
%  CAPÍTULO 7 — PRUEBAS Y EVALUACIÓN
% ============================================================
\chapter{Pruebas y evaluación}
\label{ch:pruebas}

\section{Estrategia de pruebas}
\label{sec:estrategia-pruebas}

La estrategia de pruebas del proyecto sigue una pirámide con tres niveles:

\begin{enumerate}
  \item \textbf{Pruebas unitarias:} verifican funciones aisladas de la tubería de
    transformación (mapeo de tipos, aplicación de plugins, construcción de columnas).
  \item \textbf{Pruebas de integración:} verifican el comportamiento de los endpoints
    HTTP contra una base de datos de prueba real (PostgreSQL en contenedor Docker).
  \item \textbf{Pruebas de sistema / manuales:} verificación de la interfaz de usuario
    mediante exploración manual de los flujos principales.
\end{enumerate}

\section{Pruebas del backend}
\label{sec:pruebas-backend}

\subsection{Herramientas y configuración}
\label{subsec:herramientas-test}

Las pruebas del backend utilizan \textbf{pytest}~\cite{pytest} con las siguientes
extensiones:

\begin{itemize}
  \item \texttt{pytest-asyncio}: soporte para corrutinas asíncronas en las pruebas.
  \item \texttt{httpx.AsyncClient}: cliente HTTP asíncrono para llamar a los endpoints
    FastAPI sin levantar un servidor real.
  \item \textbf{Base de datos de prueba:} instancia PostgreSQL separada, creada con
    \texttt{create\_all()} al inicio de la sesión de pruebas y destruida al final.
\end{itemize}

\subsection{Pruebas de la tubería de esquemas}
\label{subsec:pruebas-pipeline}

\begin{table}[H]
  \centering
  \caption{Casos de prueba de la tubería de transformación.}
  \label{tab:tests-pipeline}
  \begin{tabular}{lp{9cm}}
    \toprule
    \textbf{Caso} & \textbf{Descripción} \\
    \midrule
    \texttt{test\_label\_plugin} & Verifica que \texttt{x-label} se genera
      correctamente desde snake\_case. \\
    \texttt{test\_foreign\_key\_plugin} & Verifica que las claves foráneas producen
      \texttt{x-foreign-key} y \texttt{x-options-endpoint}. \\
    \texttt{test\_readonly\_plugin} & Verifica que \texttt{id} y \texttt{created\_at}
      se marcan como \texttt{x-readonly}. \\
    \texttt{test\_validator\_plugin} & Verifica que \texttt{String(255)} produce
      \texttt{maxLength: 255} en el esquema. \\
    \texttt{test\_plugin\_order} & Verifica que los plugins se aplican en orden
      de prioridad y que uno posterior puede sobrescribir a uno anterior. \\
    \bottomrule
  \end{tabular}
\end{table}

\subsection{Pruebas del endpoint bulk}
\label{subsec:pruebas-bulk}

El endpoint \texttt{/bulk} está cubierto por un conjunto de pruebas de integración que
verifican todos los escenarios relevantes:

\begin{table}[H]
  \centering
  \caption{Casos de prueba del endpoint /bulk.}
  \label{tab:tests-bulk}
  \begin{tabular}{lp{8.5cm}}
    \toprule
    \textbf{Caso} & \textbf{Descripción} \\
    \midrule
    \texttt{test\_bulk\_create\_all\_valid} & Crea múltiples filas válidas; verifica
      que todas tienen \texttt{ok: true} y reciben un \texttt{id}. \\
    \texttt{test\_bulk\_update\_existing} & Actualiza una fila existente; verifica el
      estado en BD. \\
    \texttt{test\_bulk\_delete\_op} & Elimina una fila con \texttt{\_op: "delete"};
      verifica que ya no existe en BD. \\
    \texttt{test\_bulk\_mixed} & Mezcla creación, actualización y eliminación en una
      sola petición. \\
    \texttt{test\_partial\_mode} & En modo partial, una fila inválida no bloquea las
      demás. \\
    \texttt{test\_abort\_mode} & En modo abort, un fallo causa rollback de todo el
      lote. \\
    \texttt{test\_empty\_body\_noop} & Un array vacío devuelve 200 con resumen a cero. \\
    \texttt{test\_invalid\_on\_error} & Un valor desconocido de \texttt{on\_error}
      devuelve 422. \\
    \bottomrule
  \end{tabular}
\end{table}

\subsection{Cobertura de código}
\label{subsec:cobertura}

% TODO: actualizar con los valores reales de cobertura
\todo[inline]{Ejecutar \texttt{pytest --cov=uniback} e insertar el informe de cobertura.}

La cobertura de las pruebas automatizadas del backend es del \textbf{XX\%} sobre los
módulos críticos (\texttt{crud\_factory.py}, \texttt{schema/generator.py} y los plugins).
Los módulos de configuración y arranque del servidor quedan fuera del alcance de las
pruebas automatizadas.

\section{Pruebas del frontend}
\label{sec:pruebas-frontend}

\subsection{Pruebas unitarias Angular}
\label{subsec:pruebas-angular}

Las pruebas unitarias del \textit{frontend} utilizan \textbf{Jasmine} con
\textbf{Angular TestBed}. Los módulos críticos cubiertos son:

\begin{itemize}
  \item \texttt{schema-to-formly.ts}: pruebas de mapeo de tipos y generación de
    validadores.
  \item \texttt{schema-to-aggrid.ts}: pruebas de resolución de editores y mapeo de
    columnas.
  \item \texttt{BulkEditGridComponent}: pruebas del ciclo de vida (carga de filas,
    añadir fila, marcar como eliminada, construir payload).
\end{itemize}

\subsection{Pruebas de sistema}
\label{subsec:pruebas-sistema}

Las pruebas de sistema se realizan de forma manual siguiendo guiones de prueba
predefinidos. La tabla~\ref{tab:tests-sistema} recoge los flujos principales verificados.

\begin{table}[H]
  \centering
  \caption{Guiones de prueba de sistema.}
  \label{tab:tests-sistema}
  \begin{tabular}{lll}
    \toprule
    \textbf{Flujo} & \textbf{Resultado esperado} & \textbf{Resultado} \\
    \midrule
    Navegar a browser & Tabla con datos paginados & \checkmark \\
    Crear nuevo registro & Formulario generado, guardado OK & \checkmark \\
    Editar registro existente & Campos pre-rellenos, PUT exitoso & \checkmark \\
    Eliminar registro & Confirmación + fila desaparece & \checkmark \\
    Navegar a editable\_table & Cuadrícula con filas existentes & \checkmark \\
    Añadir filas en cuadrícula & Filas en azul, contador actualizado & \checkmark \\
    Guardar en bloque (partial) & Filas válidas se guardan, inválidas en rojo & \checkmark \\
    Guardar en bloque (abort) & Fallo en una fila revierte todo & \checkmark \\
    \bottomrule
  \end{tabular}
\end{table}

\section{Evaluación del rendimiento}
\label{sec:evaluacion-rendimiento}

\todo[inline]{Incluir métricas de rendimiento: tiempo de respuesta del endpoint /bulk
con N=100, N=500, N=1000 filas; tiempo de generación de esquema para M plugins.}

Se realizaron pruebas de carga básicas sobre el endpoint \texttt{/bulk} utilizando
\textbf{Locust}~\cite{locust} con los siguientes resultados preliminares:

\begin{table}[H]
  \centering
  \caption{Rendimiento del endpoint /bulk (modo partial, PostgreSQL local).}
  \label{tab:rendimiento-bulk}
  \begin{tabular}{rrrr}
    \toprule
    \textbf{Filas} & \textbf{P50 (ms)} & \textbf{P95 (ms)} & \textbf{Errores} \\
    \midrule
    100  & -- & -- & 0\% \\
    500  & -- & -- & 0\% \\
    1000 & -- & -- & 0\% \\
    \bottomrule
  \end{tabular}
\end{table}
